\section{Orthogonal Matrices}

\textbf{(a)}
If \(U\) and \(V\) are orthogonal, then \(U^T U=I\) and \(V^T V=I\). Hence
\[
(UV)^T(UV)=V^T U^T U V=V^T V=I
\]
Thus \(UV\) has orthonormal columns. Since \(U\) and \(V\) are square, \(UV\) is square, so \(UV\) is orthogonal.

\textbf{(b)}
Since \(U\) is orthogonal,
\[
U^T U=I
\]
Therefore \(U^{-1}=U^T\), and also \(U U^T=I\). Then
\[
(U^{-1})^T U^{-1}=(U^T)^T U^T=U U^T=I
\]
Thus \(U^{-1}\) is orthogonal.

\textbf{(c)}
Let
\[
U=\begin{bmatrix}a&b\\c&d\end{bmatrix}
\]
The first column of \(U\) has norm one, so for some angle \(\theta\),
\[
\begin{bmatrix}a\\c\end{bmatrix}
=
\begin{bmatrix}\cos\theta\\ \sin\theta\end{bmatrix}
\]
The second column must be a unit vector orthogonal to the first column. In \(\mathbb{R}^2\), there are only two choices:
\[
\begin{bmatrix}-\sin\theta\\ \cos\theta\end{bmatrix},
\qquad
\begin{bmatrix}\sin\theta\\ -\cos\theta\end{bmatrix}
\]
Therefore every \(2\times 2\) orthogonal matrix has one of the two forms
\[
U=
\begin{bmatrix}
\cos\theta&-\sin\theta\\
\sin\theta& \cos\theta
\end{bmatrix}
\]
or
\[
U=
\begin{bmatrix}
\cos\theta& \sin\theta\\
\sin\theta&-\cos\theta
\end{bmatrix}
\]
The first matrix has determinant \(+1\) and is a rotation. The second matrix has determinant \(-1\) and is a reflection. Therefore, to decide which case a given orthogonal \(U\) is in, compute
\[
\det U
\]
If \(\det U=1\), it is a rotation. If \(\det U=-1\), it is a reflection.
